% =============================================================================
% Chapter: Advanced SQL
% =============================================================================

\chapter{Advanced SQL}

\begin{formulabox}[Key Concepts - MEMORIZE!]
\textbf{JOIN Types:}
\begin{center}
\begin{tabular}{lp{8cm}}
\toprule
\textbf{Type} & \textbf{Description} \\
\midrule
INNER JOIN & Μόνο matching rows \\
LEFT OUTER JOIN & Όλα από αριστερά + matching δεξιά (NULL αν δεν υπάρχει) \\
RIGHT OUTER JOIN & Όλα από δεξιά + matching αριστερά \\
FULL OUTER JOIN & Όλα από και τις δύο πλευρές \\
NATURAL JOIN & Auto-join on same column names \\
CROSS JOIN & Cartesian product (κάθε row με κάθε row) \\
\bottomrule
\end{tabular}
\end{center}
\end{formulabox}

% =============================================================================
\section{JOIN Types Visualization}

\begin{center}
\begin{tikzpicture}
    % INNER JOIN
    \begin{scope}[shift={(0,0)}]
        \fill[secondary!30] (0.6,0) circle (1);
        \fill[secondary!30] (1.4,0) circle (1);
        \begin{scope}
            \clip (0.6,0) circle (1);
            \fill[secondary] (1.4,0) circle (1);
        \end{scope}
        \node at (1,-1.8) {\textbf{INNER JOIN}};
        \node[font=\footnotesize] at (1,-2.2) {Only matching};
    \end{scope}

    % LEFT JOIN
    \begin{scope}[shift={(4,0)}]
        \fill[secondary] (0.6,0) circle (1);
        \fill[secondary!30] (1.4,0) circle (1);
        \node at (1,-1.8) {\textbf{LEFT JOIN}};
        \node[font=\footnotesize] at (1,-2.2) {All left + matching right};
    \end{scope}

    % RIGHT JOIN
    \begin{scope}[shift={(8,0)}]
        \fill[secondary!30] (0.6,0) circle (1);
        \fill[secondary] (1.4,0) circle (1);
        \node at (1,-1.8) {\textbf{RIGHT JOIN}};
        \node[font=\footnotesize] at (1,-2.2) {All right + matching left};
    \end{scope}

    % FULL OUTER JOIN
    \begin{scope}[shift={(12,0)}]
        \fill[secondary] (0.6,0) circle (1);
        \fill[secondary] (1.4,0) circle (1);
        \node at (1,-1.8) {\textbf{FULL OUTER}};
        \node[font=\footnotesize] at (1,-2.2) {All from both};
    \end{scope}
\end{tikzpicture}
\end{center}

% =============================================================================
\section{JOIN Examples}

\begin{examplebox}[JOIN Comparison]
\textbf{Tables:}

\begin{minipage}{0.45\textwidth}
\textbf{Employee:}
\begin{tabular}{ll}
\toprule
\textbf{eid} & \textbf{name} \\
\midrule
1 & Alice \\
2 & Bob \\
3 & Carol \\
\bottomrule
\end{tabular}
\end{minipage}
\begin{minipage}{0.45\textwidth}
\textbf{Department:}
\begin{tabular}{ll}
\toprule
\textbf{eid} & \textbf{dept} \\
\midrule
1 & CS \\
2 & Math \\
4 & Physics \\
\bottomrule
\end{tabular}
\end{minipage}

\vspace{0.5cm}

\textbf{INNER JOIN:}
\begin{lstlisting}[language=SQL]
SELECT * FROM Employee E INNER JOIN Department D ON E.eid = D.eid;
\end{lstlisting}
Result: \{(1, Alice, CS), (2, Bob, Math)\}

\textbf{LEFT JOIN:}
\begin{lstlisting}[language=SQL]
SELECT * FROM Employee E LEFT JOIN Department D ON E.eid = D.eid;
\end{lstlisting}
Result: \{(1, Alice, CS), (2, Bob, Math), (3, Carol, NULL)\}

\textbf{FULL OUTER JOIN:}
\begin{lstlisting}[language=SQL]
SELECT * FROM Employee E FULL OUTER JOIN Department D ON E.eid = D.eid;
\end{lstlisting}
Result: \{(1, Alice, CS), (2, Bob, Math), (3, Carol, NULL), (NULL, NULL, Physics)\}
\end{examplebox}

% =============================================================================
\section{Self-Join}

\begin{examplebox}[Self-Join: Finding Pairs]
\textbf{Query:} Find employees who work in the same department.

\begin{lstlisting}[language=SQL]
SELECT e1.name, e2.name, e1.dept
FROM Employee e1, Employee e2
WHERE e1.dept = e2.dept
  AND e1.eid < e2.eid;  -- Avoid duplicates & self-pairs
\end{lstlisting}
\end{examplebox}

% =============================================================================
\section{Subqueries}

\subsection{Types of Subqueries}

\begin{center}
\begin{tikzpicture}
    % Scalar
    \node[draw=secondary, fill=box-blue, rounded corners=5pt, minimum width=4cm, minimum height=2cm] (scalar) at (0,0) {};
    \node[anchor=north, font=\bfseries] at (scalar.north) {Scalar Subquery};
    \node[anchor=center, font=\small, align=center] at (scalar.center) {Returns single value\\WHERE x = (SELECT ...)};

    % Row
    \node[draw=success, fill=box-green, rounded corners=5pt, minimum width=4cm, minimum height=2cm] (row) at (5,0) {};
    \node[anchor=north, font=\bfseries] at (row.north) {Table Subquery};
    \node[anchor=center, font=\small, align=center] at (row.center) {Returns set of values\\WHERE x IN (SELECT ...)};

    % Correlated
    \node[draw=accent, fill=accent!20, rounded corners=5pt, minimum width=4cm, minimum height=2cm] (corr) at (10,0) {};
    \node[anchor=north, font=\bfseries] at (corr.north) {Correlated};
    \node[anchor=center, font=\small, align=center] at (corr.center) {References outer query\\WHERE EXISTS (...)};
\end{tikzpicture}
\end{center}

\subsection{EXISTS vs IN}

\begin{examplebox}[EXISTS vs IN]
\begin{lstlisting}[language=SQL]
-- Students enrolled in at least one course
-- Using IN
SELECT name FROM Student
WHERE sid IN (SELECT sid FROM Enrolls);

-- Using EXISTS (often faster for large tables)
SELECT name FROM Student S
WHERE EXISTS (SELECT 1 FROM Enrolls E WHERE E.sid = S.sid);
\end{lstlisting}

\textbf{Key Difference:}
\begin{itemize}
    \item \texttt{IN}: Builds set, then checks membership
    \item \texttt{EXISTS}: Stops at first match (short-circuit)
\end{itemize}
\end{examplebox}

% =============================================================================
\section{Window Functions}

\begin{formulabox}[Window Functions]
\begin{lstlisting}[language=SQL]
function_name() OVER (
    [PARTITION BY col]
    [ORDER BY col]
    [ROWS BETWEEN ... AND ...]
)
\end{lstlisting}

\textbf{Common Functions:}
\begin{itemize}
    \item \texttt{ROW\_NUMBER()}: Unique sequential number
    \item \texttt{RANK()}: Rank with gaps for ties
    \item \texttt{DENSE\_RANK()}: Rank without gaps
    \item \texttt{SUM(), AVG(), COUNT()}: Running aggregates
\end{itemize}
\end{formulabox}

\begin{examplebox}[ROW\_NUMBER Example]
\textbf{Query:} Top 3 highest paid employees per department.

\begin{lstlisting}[language=SQL]
WITH ranked AS (
    SELECT name, dept, salary,
           ROW_NUMBER() OVER (PARTITION BY dept
                              ORDER BY salary DESC) AS rn
    FROM Employee
)
SELECT name, dept, salary
FROM ranked
WHERE rn <= 3;
\end{lstlisting}
\end{examplebox}

% =============================================================================
\section{Common Table Expressions (CTE)}

\begin{examplebox}[CTE with WITH]
\begin{lstlisting}[language=SQL]
WITH high_earners AS (
    SELECT * FROM Employee WHERE salary > 50000
),
cs_employees AS (
    SELECT * FROM high_earners WHERE dept = 'CS'
)
SELECT * FROM cs_employees;
\end{lstlisting}
\end{examplebox}

% =============================================================================
\section{Solved Exam Problems}

\subsection{2024-25 Exam Question 7 Style: Join Algorithm}

\begin{examplebox}[Hash Join vs Merge Join vs Nested Loop]
\textbf{Question:} Which join algorithm is best for each scenario?

\begin{center}
\begin{tabular}{lp{5cm}l}
\toprule
\textbf{Algorithm} & \textbf{Best When} & \textbf{Complexity} \\
\midrule
\textbf{Nested Loop} & Small tables, or indexed inner & $O(n \cdot m)$ \\
\textbf{Hash Join} & Equality joins, one table fits in memory & $O(n + m)$ \\
\textbf{Merge Join} & Both tables sorted on join key & $O(n + m)$ \\
\bottomrule
\end{tabular}
\end{center}

\textbf{Decision Tree:}
\begin{enumerate}
    \item Sorted on join key? → \textbf{Merge Join}
    \item Equality condition + one fits memory? → \textbf{Hash Join}
    \item Small tables or index available? → \textbf{Nested Loop}
\end{enumerate}
\end{examplebox}

\subsection{Complex Query with Multiple Joins}

\begin{examplebox}[Multi-table Join]
\textbf{Tables:}
\begin{itemize}
    \item \texttt{Student(sid, name, advisor\_id)}
    \item \texttt{Professor(pid, name, dept)}
    \item \texttt{Course(cid, title, prof\_id)}
    \item \texttt{Enrolls(sid, cid, grade)}
\end{itemize}

\textbf{Query:} Find students whose advisor teaches a course they're enrolled in.

\begin{lstlisting}[language=SQL]
SELECT DISTINCT S.name
FROM Student S
JOIN Professor P ON S.advisor_id = P.pid
JOIN Course C ON P.pid = C.prof_id
JOIN Enrolls E ON S.sid = E.sid AND C.cid = E.cid;
\end{lstlisting}
\end{examplebox}

% =============================================================================
\section{Common Mistakes}

\begin{warningbox}[Συχνά Λάθη στις Εξετάσεις]
\begin{enumerate}
    \item \textbf{LEFT JOIN με WHERE condition on right table}:
    \begin{lstlisting}[language=SQL]
-- WRONG: Converts to INNER JOIN!
SELECT * FROM A LEFT JOIN B ON A.id = B.id
WHERE B.value = 'x';

-- CORRECT: Put condition in ON clause
SELECT * FROM A LEFT JOIN B ON A.id = B.id AND B.value = 'x';
    \end{lstlisting}

    \item \textbf{NATURAL JOIN surprises}:
    \begin{itemize}
        \item Matches on ALL columns with same name
        \item Can fail silently if column names change
    \end{itemize}

    \item \textbf{Correlated subquery performance}:
    \begin{itemize}
        \item Executes once per outer row
        \item Consider JOIN for better performance
    \end{itemize}
\end{enumerate}
\end{warningbox}

% =============================================================================
\section{Quick Reference}

\begin{center}
\begin{tikzpicture}
    \node[draw=ntua-blue, fill=ntua-blue!10, rounded corners=5pt, minimum width=14cm, minimum height=6cm] (main) {};

    \node[anchor=north west, font=\bfseries\Large\color{ntua-blue}] at ([xshift=5pt, yshift=-5pt]main.north west) {Advanced SQL Patterns};

    \node[anchor=north west, align=left, font=\small] at ([xshift=10pt, yshift=-35pt]main.north west) {
        \textbf{Find unmatched rows:}\\
        \texttt{SELECT * FROM A LEFT JOIN B ON ... WHERE B.id IS NULL}\\[5pt]
        \textbf{Top N per group:}\\
        \texttt{WITH r AS (SELECT *, ROW\_NUMBER() OVER (PARTITION BY ...) AS rn)}\\
        \texttt{SELECT * FROM r WHERE rn <= N}\\[5pt]
        \textbf{Running total:}\\
        \texttt{SUM(x) OVER (ORDER BY date ROWS UNBOUNDED PRECEDING)}\\[5pt]
        \textbf{Exists check:}\\
        \texttt{WHERE EXISTS (SELECT 1 FROM ... WHERE ...)}
    };

    \node[draw=secondary, fill=box-blue, rounded corners=5pt, minimum width=5cm, minimum height=1.5cm]
        at ([xshift=-4cm, yshift=1.3cm]main.south) {};
    \node[align=center, font=\footnotesize] at ([xshift=-4cm, yshift=1.3cm]main.south) {
        \textbf{INNER}: Only matches\\
        \textbf{LEFT}: All left + matches
    };

    \node[draw=accent, fill=box-red, rounded corners=5pt, minimum width=5cm, minimum height=1.5cm]
        at ([xshift=2cm, yshift=1.3cm]main.south) {};
    \node[align=center, font=\footnotesize] at ([xshift=2cm, yshift=1.3cm]main.south) {
        \textbf{Hash Join}: Equality, memory\\
        \textbf{Merge Join}: Sorted data
    };
\end{tikzpicture}
\end{center}
